% preamble-exam-zh.tex — 共享 preamble
% 用于所有 exam-zh 模板
%
% 样式策略：
%   屏幕 → 语义彩色（蓝=关键想法，红=易错，绿=路线，灰=提示/教师备注）
%   打印 → 自动降级为黑白（通过 print mode 或 monochrome 选项）

\documentclass{exam-zh}

% ---------- 基础配置 ----------
\examsetup{
  page/size = a4paper,
  page/show-head = false,
  page/show-foot = false,
  sealline/show = false,
  font = times,
  math-font = xits,
}

\usepackage{tcolorbox}
\usepackage{needspace}
\tcbuselibrary{skins,breakable}

% ---------- 打印/屏幕颜色切换 ----------
% 用法：\documentclass[monochrome]{exam-zh} 即可切换为纯黑白
% 注意：不要使用 \if@xxx 放在 makeatother 外部；这里改成普通条件名。
\newif\ifeduprintcolor
\eduprintcolortrue

\makeatletter
\@ifclasswith{exam-zh}{monochrome}{\eduprintcolorfalse}{}
\makeatother

% ---------- 语义颜色定义 ----------
% 屏幕：有色彩区分；打印：降级为灰度
\ifeduprintcolor
  % --- 屏幕彩色模式 ---
  \definecolor{edu-blue}{RGB}{47, 82, 122}       % 关键想法
  \definecolor{edu-blue-bg}{RGB}{243, 246, 252}
  \definecolor{edu-red}{RGB}{180, 40, 40}         % 易错提醒
  \definecolor{edu-red-bg}{RGB}{255, 245, 240}
  \definecolor{edu-green}{RGB}{34, 100, 60}       % 路线图
  \definecolor{edu-green-bg}{RGB}{242, 250, 245}
  \definecolor{edu-orange}{RGB}{160, 90, 0}       % 训练目标
  \definecolor{edu-orange-bg}{RGB}{255, 248, 235}
  \definecolor{edu-gray}{RGB}{100, 100, 100}      % 提示/教师备注
  \definecolor{edu-gray-bg}{RGB}{248, 248, 248}
  \definecolor{edu-purple}{RGB}{90, 50, 120}      % 思考题
  \definecolor{edu-purple-bg}{RGB}{248, 244, 252}
\else
  % --- 打印黑白模式 ---
  \definecolor{edu-blue}{gray}{0.15}
  \definecolor{edu-blue-bg}{gray}{0.96}
  \definecolor{edu-red}{gray}{0.25}
  \definecolor{edu-red-bg}{gray}{0.97}
  \definecolor{edu-green}{gray}{0.20}
  \definecolor{edu-green-bg}{gray}{0.97}
  \definecolor{edu-orange}{gray}{0.25}
  \definecolor{edu-orange-bg}{gray}{0.98}
  \definecolor{edu-gray}{gray}{0.40}
  \definecolor{edu-gray-bg}{gray}{0.97}
  \definecolor{edu-purple}{gray}{0.30}
  \definecolor{edu-purple-bg}{gray}{0.98}
\fi

% ---------- 自定义教学环境 ----------
% 共性：enhanced + breakable（跨页不断裂）

% 原题卡片 — 强边框，突出显示
\newtcolorbox{problemcard}{
  enhanced, breakable,
  colback=edu-blue-bg,
  colframe=edu-blue,
  boxrule=1.2pt,
  arc=0pt,
  left=3mm, right=3mm, top=3mm, bottom=3mm,
}

% 解题路线图 — 左边线 + 浅底
\newtcolorbox{routebox}{
  enhanced, breakable,
  colback=edu-green-bg,
  colframe=edu-green!30,
  boxrule=0pt,
  borderline west={2.5pt}{0pt}{edu-green},
  left=3mm, right=3mm, top=3mm, bottom=3mm,
}

% 关键想法 — 蓝色左边线
\newtcolorbox{keyideabox}{
  enhanced, breakable,
  colback=edu-blue-bg,
  colframe=edu-blue!30,
  boxrule=0pt,
  borderline west={2.5pt}{0pt}{edu-blue},
  left=3mm, right=3mm, top=3mm, bottom=3mm,
}

% 易错提醒 — 红色左边线
\newtcolorbox{mistakebox}{
  enhanced, breakable,
  colback=edu-red-bg,
  colframe=edu-red!20,
  boxrule=0pt,
  borderline west={3pt}{0pt}{edu-red},
  left=3mm, right=3mm, top=2.5mm, bottom=2.5mm,
}

% 提示 — 灰色虚线左边线
\newtcolorbox{hintbox}{
  enhanced, breakable,
  colback=edu-gray-bg,
  colframe=edu-gray!20,
  boxrule=0pt,
  borderline west={2pt}{0pt}{edu-gray, dashed},
  left=3mm, right=3mm, top=2mm, bottom=2mm,
}

% 思考题 — 紫色虚线左边线
\newtcolorbox{thinkbox}{
  enhanced, breakable,
  colback=edu-purple-bg,
  colframe=edu-purple!20,
  boxrule=0pt,
  borderline west={2pt}{0pt}{edu-purple, dashed},
  left=2.5mm, right=2.5mm, top=2.5mm, bottom=2.5mm,
}

% 教师备注 — 灰色左边线，小字
\newtcolorbox{teachernote}{
  enhanced, breakable,
  colback=edu-gray-bg,
  colframe=edu-gray!15,
  boxrule=0pt,
  borderline west={2pt}{0pt}{edu-gray!60},
  left=2.5mm, right=2.5mm, top=2.5mm, bottom=2.5mm,
  fontupper=\small,
  colupper=black!60,
}

% 训练目标 — 橙色左边线，小字
\newtcolorbox{traininggoal}{
  enhanced, breakable,
  colback=edu-orange-bg,
  colframe=edu-orange!15,
  boxrule=0pt,
  borderline west={1.5pt}{0pt}{edu-orange},
  left=2mm, right=2mm, top=1mm, bottom=1mm,
  fontupper=\small,
}

% 答题区 — 无边框，纯留白
\newcommand{\answerarea}[1][40mm]{%
  \vspace{2mm}%
  \begin{tcolorbox}[
    enhanced, breakable,
    colback=white,
    colframe=white,
    boxrule=0pt,
    height=#1,
    valign=top,
    bottom=0mm,
    after skip=0mm,
  ]
  \end{tcolorbox}%
}

% ---------- 字体设置 ----------
% exam-zh (ctexbook) already sets CJK fonts via ctex.
% Only override if specific fonts are needed.
% macOS: Songti SC, STHeiti, STFangsong
% Linux: SimSun, SimHei, FangSong

% ---------- 页面设置 ----------
% exam-zh (ctexbook) already loads geometry.
% Override margins if needed:
% \geometry{top=16mm, bottom=16mm, left=14mm, right=14mm}

\examsetup{
  question/show-answer = true,
  fillin/show-answer = true,
  solution/show-solution = show-stay,
}

% ---------- 教师版解答样式 ----------
\newtcolorbox{solutionblock}{
  enhanced, breakable,
  colback=edu-blue-bg,
  colframe=edu-blue!25,
  boxrule=0pt,
  borderline west={2pt}{0pt}{edu-blue!50},
  arc=0pt,
  left=3mm, right=3mm, top=2mm, bottom=2mm,
  before skip=2mm,
  after skip=3mm,
  fontupper=\small,
}

% 解答步骤编号
\newcounter{solstep}
\newcommand{\solstep}[2]{%
  \stepcounter{solstep}%
  \par\medskip
  \noindent\hangindent=6mm
  {\color{edu-blue}\bfseries\textcircled{\scriptsize\thesolstep}\enspace #1}\enspace #2\par
}

\begin{document}

\section*{三垂直与旋转等腰直角专题（三） · 练习卷（教师版）}

\section*{一、选择题（每题 12 分，共 48 分）}


\begin{keyideabox}
  \textbf{中等思想：多解与退化情形}
  在没有给出具体象限限制时，作等腰直角三角形第三顶点一般有\textbf{两个对称的解}（由于旋转方向可以顺时针或逆时针）。必须逐一讨论，即使垂足重合也不应漏解。
\end{keyideabox}



\begin{question}[points=12]
  直线 $AB$ 经过 $A(-1, 1)$ 和 $B(2, 0)$ 两点，且与 $y$ 轴交于点 $C$。若点 $D(0, n)$ 在点 $C$ 的上方，且 $\triangle ABD$ 的面积为 2，则 $n$ 的值为
  \begin{choices}
    \item 2
    \item 3
    \item 4
    \item 1.5
  \end{choices}
  \paren[A]
  \begin{solutionblock}
    直线 $AB: y = -\frac{1}{3}x + \frac{2}{3} \implies C(0, \frac{2}{3})$。\\ $D(0,n)$ 位于 $y$ 轴，$CD = n - \frac{2}{3}$。以 $CD$ 为底， $A, B$ 横坐标绝对值之和为高之和 $= 1+2=3$。\\ 面积 $S = \frac{1}{2} \times CD \times 3 = \frac{3}{2}(n-\frac{2}{3}) = 2 \implies n-\frac{2}{3} = \frac{4}{3} \implies n=2$。
  \end{solutionblock}
\end{question}



\begin{question}[points=12]
  已知 $B(2, 0)$， $D(0, 2)$。以 $BD$ 为直角边作等腰 Rt$\triangle BDP$ （直角顶点为 $B$），则满足条件的点 $P$ 坐标为
  \begin{choices}
    \item $(4, 2)$ 或 $(0, -2)$
    \item $(4, 2)$
    \item $(0, -2)$
    \item $(4, 4)$
  \end{choices}
  \paren[A]
  \begin{solutionblock}
    直角在 $B$，利用一线三垂直。过 $B(2,0)$ 作 $x$ 轴辅助线。$D(0,2) \implies D_1(0,0)$。$BD_1 = 2, DD_1=2$。\\ 全等 $\triangle BDD_1 \cong \triangle PP_1B \implies PP_1 = 2, BP_1=2$。\\ 若 $P_1$ 在 $B$ 右侧：$x_p = 4, y_p = 2$ $\implies P(4, 2)$。\\ 若 $P_1$ 在 $B$ 左侧：$x_p = 0, y_p = -2$ $\implies P(0, -2)$。所以有两个解。
  \end{solutionblock}
\end{question}



\begin{question}[points=12]
  平面直角坐标系内两点 $A(1, 1)$ 和 $B(3, 2)$。若以 $A$ 为直角顶点在直线上方作等腰 Rt$\triangle ABP$，则点 $P$ 的横坐标为
  \begin{choices}
    \item 0
    \item 2
    \item 0 或 2
    \item 1
  \end{choices}
  \paren[A]
  \begin{solutionblock}
    直角在 $A(1,1)$。过 $A$ 水平线。$B(3,2)$ 垂足 $B_1(3,1)$。$AB_1 = 2, BB_1=1$。\\ 全等 $\triangle ABB_1 \cong \triangle PP_1A \implies PP_1 = 2, AP_1=1$。\\ 在直线上方，则 $P_1$ 在 $A$ 左侧 $1$ 个单位 $\implies x_p = 0$。
  \end{solutionblock}
\end{question}



\begin{question}[points=12]
  已知点 $A(-1, 1)$。则点 $A$ 到 $y$ 轴的距离为
  \begin{choices}
    \item 1
    \item -1
    \item 2
    \item 0
  \end{choices}
  \paren[A]
  \begin{solutionblock}
    到 $y$ 轴距离为横坐标的绝对值，即 $|-1| = 1$。
  \end{solutionblock}
\end{question}


\section*{二、填空题（每题 12 分，共 12 分）}


\begin{question}[points=12]
  在平面直角坐标系中，已知 $A(0, 3)$ 和 $B(2, 1)$。以 $B$ 为直角顶点作等腰 Rt$\triangle ABP$，若点 $P$ 在 $x$ 轴上方，则点 $P$ 的坐标为 \fillin。
  \paren[$(4, 3)$]
  \begin{solutionblock}
    直角顶点在 $B(2,1)$。过 $B$ 作水平辅助线。$A(0,3)$ 垂足 $A_1(0,1)$。$BA_1 = 2, AA_1 = 2$。\\ 全等 $\triangle BAA_1 \cong \triangle PP_1B \implies PP_1 = 2, BP_1 = 2$。\\ 在 $x$ 轴上方，则 $y_p = 1+2=3$。$P_1$ 在 $B$ 右侧 $\implies x_p = 2+2=4$。坐标为 $(4, 3)$。
  \end{solutionblock}
\end{question}


\section*{三、解答题（共 40 分）}

\needspace{8\baselineskip}

\begin{problem}[points=20]
  直线 $AB$ 经过 $A(-1, 1)$ 和 $B(2, 0)$ 两点，且与 $y$ 轴交于点 $C$。点 $D(0, n)$ 位于点 $C$ 的上方。

(1) 若 $\triangle ABD$ 的面积为 2，求点 $D$ 的坐标；
(2) 在(1)的条件下，以 $BD$ 为直角边作等腰 Rt$\triangle BDP$ （直角顶点为 $B$，$\angle DBP = 90^\circ$），求满足条件的点 $P$ 的坐标。
  \answerarea[65mm]
  \setcounter{solstep}{0}
  \begin{solutionblock}
    \textbf{答案：}(1) $D(0, 2)$；(2) $P(4, 2)$ 或 $P(0, -2)$\par\medskip
    \solstep{求解析式与 D 坐标}{直线 $AB$ 解析式为 $y = -\frac{1}{3}x + \frac{2}{3} \implies C\left(0, \frac{2}{3}\right)$。\\ $D(0,n)$ 位于 $y$ 轴，$CD = n - \frac{2}{3}$。\\ $S_{\triangle ABD} = S_{\triangle ACD} + S_{\triangle BCD} = \frac{1}{2} CD \times 1 + \frac{1}{2} CD \times 2 = \frac{3}{2} CD$。\\ 令 $\frac{3}{2}\left(n - \frac{2}{3}\right) = 2 \implies n - \frac{2}{3} = \frac{4}{3} \implies n = 2$。故 $D(0, 2)$。}
    \solstep{一线三垂直求 P 坐标情况一}{直角在 $B(2,0)$。过 $B$ 做水平辅助线（即 $x$ 轴）。\\ $D(0,2)$ 垂足为 $D_1(0,0)$。垂线段 $DD_1 = 2$，$BD_1 = 2$。\\ 设另一点 $P$ 垂足为 $P_1$。由全等 $\triangle BDD_1 \cong \triangle PP_1B$ (AAS) 得：\\ $PP_1 = BD_1 = 2$，$BP_1 = DD_1 = 2$。\\ 情况一：若 $P_1$ 在 $B$ 右侧，则 $x_p = 2+2=4$。且 $P$ 在 $x$ 轴上方 $\implies y_p = 2$。\\ 坐标为 $P_1(4, 2)$。}
    \solstep{情况二与检验}{情况二：若 $P_1$ 在 $B$ 左侧，则 $x_p = 2-2=0$。且 $P$ 在 $x$ 轴下方 $\implies y_p = -2$。\\ 坐标为 $P_2(0, -2)$。\\ 这两个点均满足条件，且 $P_2(0,-2)$ 的垂足正好是 $D_1(0,0)$ (退化解)。\\ 故 $P(4, 2)$ 或 $P(0, -2)$。}
  \end{solutionblock}
\end{problem}


\needspace{8\baselineskip}

\begin{problem}[points=20]
  已知平面直角坐标系中两点 $A(0, 3)$ 和 $B(2, 1)$。

(1) 若以 $AB$ 为直角边作等腰 Rt$\triangle ABP$ （直角顶点为 $B$，$\angle ABP = 90^\circ$），求出所有满足条件的点 $P$ 坐标；
(2) 在坐标轴上画出对应的示意图，并写出全等三角形的对应边关系。
  \answerarea[60mm]
  \setcounter{solstep}{0}
  \begin{solutionblock}
    \textbf{答案：}(1) $P(4, 3)$ 或 $P(0, -1)$；(2) $PP_1 = BA_1 = 2$, $BP_1 = AA_1 = 2$\par\medskip
    \solstep{构建辅助线与全等}{过直角顶点 $B(2,1)$ 做水平辅助线 $y=1$。\\ $A(0,3)$ 垂足 $A_1(0,1)$。$BA_1 = 2, AA_1 = 2$。\\ 设点 $P$ 垂足为 $P_1$。由 AAS 得 $\triangle BAA_1 \cong \triangle PP_1B$。\\ 对应边：$PP_1 = BA_1 = 2$，$BP_1 = AA_1 = 2$。}
    \solstep{求解情况一与情况二}{情况一：$P$ 位于辅助线上方，且 $P_1$ 在 $B$ 右侧。\\ $y_p = 1+2=3$；$x_p = 2+2=4$。坐标为 $P_1(4, 3)$。\\ 情况二：$P$ 位于辅助线下方，且 $P_1$ 在 $B$ 左侧。\\ $y_p = 1-2=-1$；$x_p = 2-2=0$。坐标为 $P_2(0, -1)$。\\ 故点 $P$ 坐标为 $(4, 3)$ 或 $(0, -1)$。}
  \end{solutionblock}
\end{problem}


\clearpage
\section*{参考答案}


\begin{keyideabox}
  \textbf{选择题}
  1. A  2. A  3. A  4. A
\end{keyideabox}



\begin{keyideabox}
  \textbf{填空题}
  1. (4, 3)
\end{keyideabox}



\begin{keyideabox}
  \textbf{解答题}
  第 1 题：(1) $D(0, 2)$；(2) $P(4, 2)$ 或 $P(0, -2)$
第 2 题：(1) $P(4, 3)$ 或 $P(0, -1)$；(2) $PP_1 = BA_1 = 2$, $BP_1 = AA_1 = 2$
\end{keyideabox}



\end{document}
